\chapter{Refactorizaciones}

\section{Eliminar código muerto y duplicado de la interfaz}
\begin{itemize}
\item \textbf{\texttt{MenuController} eliminado.} No tenía ninguna referencia (ni en FXML ni en Java): la aplicación navega directamente a la vista del rol autenticado.
\item \textbf{\texttt{AyudasUI} (nuevo).} Centraliza \texttt{animarNavegacion(Button)} — las tres copias idénticas en \texttt{AdminController}, \texttt{VendedorController} y \texttt{AlmacenController} — y \texttt{animarHover(Button)}, la variante leve del login. Se eliminaron las cuatro copias locales.
\item \textbf{\texttt{InformeTexto} (nuevo).} Unifica la generación del texto plano de informes (\texttt{borde()}, \texttt{fila()}, \texttt{ajustar()}) que estaba duplicada en \texttt{AdminController} (informes de ventas y entradas) y \texttt{VendedorController} (proforma). Consecuencia estética menor: las celdas que exceden el ancho se truncan con «…» también en la proforma, como ya hacían los informes.
\item \textbf{Referencias directas.} Se sustituyó el prefijo cualificado \texttt{com.overbyte.Navegador.} (~60 usos) por \texttt{Navegador} importado, y se eliminaron los wrappers locales \texttt{mostrarAlerta()} que solo delegaban en \texttt{Navegador.avisar()}.
\end{itemize}

\section{Centralizar nombres de archivos CSV}
Nueva clase \texttt{config/ArchivosCSV} con los nombres canónicos de los nueve CSV que usa el sistema (\texttt{INVENTARIO}, \texttt{VENTAS}, \texttt{CLIENTES}, \texttt{GARANTIAS}, \texttt{CAJA}, \texttt{EMPLEADOS}, \texttt{PROVEEDORES}, \texttt{ENTRADAS}, \texttt{PEDIDOS}) y con \texttt{INVENTARIO\_ENCABEZADO}, la cabecera que estaba duplicada en \texttt{Gestor} y \texttt{Almacen}. Incluye el reporte de \texttt{MotorCSV} y las lecturas de los tres controladores. No queda ningún literal de nombre de archivo fuera de esa clase.

\section{Centralizar el guardado de archivos de texto}
Nueva clase \texttt{AyudasArchivo} con dos utilidades compartidas:
\begin{itemize}
\item \texttt{elegirRutaTxt(ventana, titulo, nombreInicial)}: \texttt{FileChooser} de guardado con extensión \texttt{.txt} garantizada (o \texttt{null} si el usuario cancela).
\item \texttt{escribirTxt(ruta, contenido)}: escritura UTF-8 con manejo de \texttt{IOException}; devuelve \texttt{boolean} para que cada controlador conserve su mensaje de éxito o error específico.
\end{itemize}
Sustituye cuatro bloques duplicados: guardar venta confirmada, guardar proforma y exportar informe de ventas e informe de entradas.

\section{Unificar la validación de la consola}
\texttt{Main.java} (punto de entrada de consola) eliminó sus validadores locales \texttt{textoValido()} y \texttt{datosClienteValidos()} — idénticos en intención pero más estrictos que \texttt{ValidadorDatos} — y delega en \texttt{services.ValidadorDatos}, la misma fuente de verdad que usa la GUI. Consecuencia: la consola ahora acepta nombres alfanuméricos (p. ej. \emph{Juan123}) como siempre hizo la GUI. Se ajustaron cuatro pruebas de \texttt{MainTest} (\texttt{nombreEmpleadoSoloNumerosEsRechazado}, \texttt{proveedorNombreSoloNumeros}, \texttt{altaEmpleadoRechazadaNoPersiste}, \texttt{telefonoSinDigitosSuficientes}): el caso de rechazo queda cubierto por nombres compuestos solo de dígitos, que ambas reglas continúan rechazando.

\section{El enum Categoria}
Se introdujo en \texttt{2fe216f} como fuente de verdad para el filtro, y se retiró en \texttt{2ecf69d} porque fijar las categorías en código obliga a recompilar ante cualquier categoría nueva del establecimiento. La funcionalidad no se perdió: el filtro autodescubre lo ya listado en \texttt{data/inventario.csv} (ver Capítulo~\ref{chap:cambios}).
